\section*{الفصل \ref{ch:b3:advanced-geochronology} --- الجيوكرونولوجيا المتقدّمة}

\begin{solution}{exo:b3:advanced-geochronology:1}
$t_{206} = \ln(1.04500)/\lambda_{238} = \qty{283.8}{\Ma}$; $t_{207} =
\ln(1.3300)/\lambda_{235} = \qty{289.6}{\Ma}$. تختلف الأعمار بمقدار
\qty{6}{\Ma}: الزركون غير متوافق. و$t_{207}$ هو الأقدم، لذا يقع
أسفل \omterm{def:b3:advanced-geochronology:concordia}{الكونكورديا}، على جانب فقد الرصاص.
\end{solution}

\begin{solution}{exo:b3:advanced-geochronology:2}
$J = (e^{5.543\times10^{-10}\times2.8201\times10^7} - 1)/1.3129 =
0.015754/1.3129 = 0.01200$. ثم $t = \ln(1 + 0.01200\times10.000)/\lambda =
\ln 1.12/(5.543\times10^{-10}) = \qty{204.5}{\Ma}$.
\end{solution}

\begin{solution}{exo:b3:advanced-geochronology:3}
$T_{\mathrm{CHUR}} = \ln\bigl(1 + (0.51180 - 0.512638)/(0.1200 -
0.1966)\bigr)/\lambda = \ln(1.010940)/\lambda = \qty{1.67}{\Ga}$;
$T_{\mathrm{DM}} = \ln\bigl(1 + (0.51180 - 0.51315)/(0.1200 - 0.2137)\bigr)/\lambda
= \ln(1.014408)/\lambda = \qty{2.19}{\Ga}$.
\end{solution}

\begin{solution}{exo:b3:advanced-geochronology:4}
بدون تفكّك، $t = N/P_0 = \num{2.0e5}/4.96 = \qty{40.3}{kyr}$. مع التفكّك،
$t = -\ln(1 - N\lambda/P_0)/\lambda = -\ln(1 - 0.02014)/\lambda =
\qty{40.7}{kyr}$، أي أقدم بنسبة \qty{1}{\percent}.
\end{solution}

\begin{solution}{exo:b3:advanced-geochronology:5}
من أجل $t$ صغيرة، $x \approx \lambda_{235}t$ و$y \approx \lambda_{238}t$، ومنه
$y \approx (\lambda_{238}/\lambda_{235})x$، وميله $0.1575$. إذا فُقد الرصاص
حديثًا جدًّا، $t_2 \approx 0$ ويكون \omterm{def:b3:advanced-geochronology:concordia}{التقاطع السفلي} $C(0)$، أي
نقطة الأصل. على مستقيم يمرّ بنقطة الأصل يكون $y/x$ ثابتًا، وكذلك
$^{207}\mathrm{Pb}^*/^{206}\mathrm{Pb}^* = x/(137.8\,y)$: كلّ بلّورة تعطي
عمر الرصاص--الرصاص للتقاطع العلوي، أي عمر التكوين.
\end{solution}

\begin{solution}{exo:b3:advanced-geochronology:6}
الأوزان $1.5625$، $1$، $1.2346$، $2.0408$، $W = 5.838$؛ $\bar x =
\qty{296.09}{\Ma}$، $\sigma = 1/\sqrt W = \qty{0.41}{\Ma}$. مجموع مربعات
الانحرافات الموزونة يساوي $2.32$؛ مع $f = 3$، $\mathrm{MSWD} = 0.77$، أدنى من
$1 + 2\sqrt{2/3} = 2.63$: الأعمار متوافقة.
\end{solution}

\begin{solution}{exo:b3:advanced-geochronology:7}
$\lambda\zeta g\rho_d\rho_s/\rho_i = 1.55125\times10^{-10}\times350\times0.5
\times\num{1.10e6}\times0.5 = 0.014931$; $t = \ln(1.014931)/\lambda =
\qty{95.5}{\Ma}$. إذا بَرد الصخر بسرعة، فإنّه يؤرّخ مروره عبر
نطاق التلدين الجزئي للأباتيت، نحو 60 إلى \qty{120}{\celsius}؛ وإذا
بَرد ببطء، فهو عمر مختلط، ويجب أن تُبيّن أطوال الآثار ذلك.
\end{solution}

\begin{solution}{exo:b3:advanced-geochronology:8}
لكلّ غرام: اليورانيوم $20\times10^{-6}/238.05\times N_A = \num{5.06e16}$ ذرّة، منها
\num{5.02e16} من اليورانيوم-238 و\num{3.64e14} من اليورانيوم-235؛ الثوريوم
\num{1.038e17}، $\lambda_{232} = \ln 2/\num{1.40e10} = \qty{4.95e-11}{yr^{-1}}$;
الهليوم $10^{-9}N_A = \num{6.02e14}$. \omterm{def:b3:advanced-geochronology:cosmogenic}{معدّل الإنتاج} هو
$8\lambda_{238}{}^{238}\mathrm U + 7\lambda_{235}{}^{235}\mathrm U +
6\lambda_{232}\mathrm{Th} = (6.23 + 0.25 + 3.08)\times10^7 =
\num{9.57e7}\text{ ذرّة}\,\unit{g^{-1}.yr^{-1}}$، و$t = \num{6.02e14}/\num{9.57e7} =
\qty{6.29}{Myr}$. بعد التصحيح: $6.29/0.75 = \qty{8.39}{\Ma}$.
\end{solution}

\begin{solution}{exo:b3:advanced-geochronology:9}
$P_0/\lambda = 4.96/\num{4.994e-7} = \num{9.93e6}\text{ ذرّة}\,\unit{g^{-1}}$، ويبلغ
\qty{90}{\percent} بعد $\ln 10/\lambda = \qty{4.6}{Myr}$. عند
\qty{10}{m.Myr^{-1}} $= \qty{1e-3}{cm.yr^{-1}}$: $\rho\varepsilon/\Lambda =
2.7\times10^{-3}/160 = \qty{1.69e-5}{yr^{-1}}$, $\mu = \qty{1.74e-5}{yr^{-1}}$
و$N = 4.96/\num{1.74e-5} = \num{2.85e5}\text{ ذرّة}\,\unit{g^{-1}}$، أي أقلّ بنحو 35 مرّة.
\end{solution}

\begin{solution}{exo:b3:advanced-geochronology:10}
مع $W = w_1 + w_2$، $x_1 - \bar x = w_2(x_1 - x_2)/W$ و$x_2 - \bar x =
-w_1(x_1 - x_2)/W$، ومنه $\sum w_i(x_i - \bar x)^2 = (x_1 - x_2)^2(w_1w_2^2 +
w_2w_1^2)/W^2 = (x_1 - x_2)^2w_1w_2/W = (x_1 - x_2)^2/(\sigma_1^2 +
\sigma_2^2)$؛ مع $f = 1$ يكون هذا هو MSWD. إذا قاس الاثنان العمر نفسه،
$x_1 - x_2$ له متوسط 0 وتباين $\sigma_1^2 + \sigma_2^2$، ومنه فإنّ
يكون توقّع مربّعه مقسومًا على هذا التباين يساوي 1.
\end{solution}

\begin{solution}{exo:b3:advanced-geochronology:11}
$C(2700) = (13.283, 0.52019)$، $C(500) = (0.63628, 0.08065)$. $f = (0.52019 -
0.300)/(0.52019 - 0.08065) = 0.501$، نصف رصاصه. $x = 0.499\times13.283 +
0.501\times0.63628 = 6.948$. $t_{206} = \qty{1691}{\Ma}$، $t_{207} =
\qty{2105}{\Ma}$؛ النسبة $x/(137.8\,y) = 0.1681$ تعطي بالتكرار عمرًا
عمر الرصاص--الرصاص \qty{2537}{\Ma}. لا معنى لأيّ من الأعمار الثلاثة، فهي
خلائط من الحدثين، ولا يؤرّخها إلا التقاطعات. عمر الرصاص--الرصاص
هو حدّ أدنى لزمن التبلّور.
\end{solution}

\begin{solution}{exo:b3:advanced-geochronology:12}
بدون تعرية: $t = -\ln(1 - \num{1.0e6}\times\num{4.994e-7}/4.96)/\lambda =
\qty{212}{kyr}$. الحالة المستقرّة: $\rho\varepsilon/\Lambda = P_0/N - \lambda =
\qty{4.46e-6}{yr^{-1}}$، $\varepsilon = \qty{2.64e-4}{cm.yr^{-1}} =
\qty{2.6}{m.Myr^{-1}}$. لأيّ $\varepsilon$ أدنى من هذه القيمة، $N =
P_0(1 - e^{-\mu t})/\mu$ يحدّد زمنًا واحدًا $t$، يتراوح من \qty{212}{kyr} صعودًا حتى
اللانهاية: لدى القياس الواحد مجهولان. تفصلهما نويدة ثانية ذات
عمر نصف آخر، مثل الألومنيوم-26 في الكوارتز نفسه، أو تركيز النويدة عند
عمق في لُبّ الرواسب.
\end{solution}

\begin{solution}{pb:b3:advanced-geochronology:1}
\textbf{1.} $t_{206}$ / $t_{207}$، بوحدة \unit{\Ma}: \babelsublr{Z1 301.9 / 302.2}، \babelsublr{Z2 284.1 / 286.5}،
\babelsublr{Z3 265.7 / 268.5}، \babelsublr{Z4 246.1 / 250.1}، \babelsublr{Z5 225.1 / 229.4}. خطأ قدره
\qty{0.2}{\percent} في $y$ نحو \qty{0.6}{\Ma}: الزركونة Z1 متوافقة ضمن
الخطأ، والزركونات الأخرى غير متوافقة، ويتزايد عدم التوافق من Z2 إلى Z5.

\textbf{2.} فقدت الزركونات كسورًا متزايدة من رصاصها في حدث
لاحق واحد: فهي تقع أبعد فأبعد أسفل الوتر. ويقع الوتر أسفل
\omterm{def:b3:advanced-geochronology:concordia}{الكونكورديا} المقعّرة، حيث يكون $x$ أكبر لقيمة معلومة من $y$: $t_{207} > t_{206}$.

\textbf{3.} الزركونة \mbox{\textdir TLT Z1}: $X = 1/0.04794 = 20.859$، $Y = 0.3467/(137.8\times0.04794) =
0.05248$. الزركونة \mbox{\textdir TLT Z5}: $X = 28.145$، $Y = 0.05178$.

\textbf{4.} الأوزان $1.000$، $1.235$، $1.563$، $1.563$، $2.041$ (بوحدة
$10^8$)؛ $\bar x = 0.29402$، $\bar y = 0.040892$، $S_{xx} = \num{7.845e5}$.

\textbf{5.} $b = 0.13260$، $a = \bar y - b\bar x = 0.001905$.

\textbf{6.} عند \qty{310}{\Ma}، $\mathrm dy/\mathrm dx = 0.1575\,
e^{(\lambda_{238} - \lambda_{235})t} = 0.1218$. \omterm{def:b3:advanced-geochronology:concordia}{الديسكورديا} أشدّ انحدارًا من \omterm{def:b3:advanced-geochronology:concordia}{الكونكورديا}
بنسبة \qty{9}{\percent} فقط: تلتقي معها بزاوية صغيرة،
وسيكون \omterm{def:b3:advanced-geochronology:concordia}{التقاطع العلوي} غير دقيق.

\textbf{7.} بطريقة نيوتن ابتداءً من \qty{310}{\Ma}: $t = \qty{311.2}{\Ma}$.

\textbf{8.} ابتداءً من \qty{100}{\Ma}: $t = \qty{99.7}{\Ma}$.

\textbf{9.} $(y_i - a - bx_i)/\sigma_i = 0.64$، $-0.91$، $0.35$، $-0.24$،
$0.17$؛ $\mathrm{MSWD} = 1.435/3 = 0.48$، أقل من الحدّ $2.63$ عند $f = 3$:
تتوافق الزركونات مع خطّ واحد.

\textbf{10.} $s_b = 1/\sqrt{\num{7.845e5}} = 0.00113$.

\textbf{11.} عند $b - 2s_b = 0.13034$ يكون \omterm{def:b3:advanced-geochronology:concordia}{التقاطع العلوي}
\qty{322.9}{\Ma}؛ وعند $b + 2s_b = 0.13485$ يكون \qty{302.6}{\Ma}. ويكون العمر
$311 \pm \qty{10}{\Ma}$ عند انحرافين معياريين.

\textbf{12.} $f = (x(311.2) - x_i)/(x(311.2) - x(99.7))$: الزركونة \mbox{\textdir TLT Z1} \qty{4.7}{\percent}،
\mbox{\textdir TLT Z5} \qty{41}{\percent}.

\textbf{13.} $J = 0.015754/1.3129 = 0.0119999$، $0.01200$.

\textbf{14.} الخطوة 1: \qty{150.0}{\Ma}؛ الخطوة 4: \qty{296.6}{\Ma}؛ الخطوة 10:
\qty{295.0}{\Ma}.

\textbf{15.} الخطوات من 4 إلى 10: سبع خطوات متتالية، تمثّل \qty{82}{\percent} من الأرغون-39،
دون اتجاه واضح. وتتزايد الخطوات من 1 إلى 3 على شكل درج.

\textbf{16.} أخطاء متساوية: المتوسط هو $2072.9/7 = \qty{296.1}{\Ma}$، $\sigma =
1.1/\sqrt7 = \qty{0.42}{\Ma}$؛ ومجموع مربعات الانحرافات يساوي $3.83$،
$\mathrm{MSWD} = 3.83/(1.21\times6) = 0.53$، أقل من $2.15$.

\textbf{17.} $\bar R = \sum f_iR_i = 14.270$، $t = \qty{285.2}{\Ma}$: إنّ
الأرغون المفقود من الأطراف، المحرَّر في الخطوات الأولى، يخفّض المتوسط،
كما يخفّض عمر البوتاسيوم--الأرغون.

\textbf{18.} $(700 - 375)/(311.2 - 296.1) = \qty{21.5}{\celsius.Myr^{-1}}$
(مع مدى واسع، لأنّ عمر الزركون معلوم بدقّة $\pm\qty{10}{\Ma}$).

\textbf{19.} قرب \qty{100}{\Ma} كان الغرانيت مدفونًا تحت
غطائه الميزوزوي، وأدّى التسخين أو دوران السوائل إلى فقدان الزركونات للرصاص.
فقد البيوتيت الأرغون من أطرافه، وهو ما سجّله الدرج؛ وأدنى
خطوة له، \qty{150}{\Ma}، هي خليط من أرغون قديم وأرغون أُعيد ضبطه. ولم تتأثّر
هضبة العمر: لم يُسخَّن البيوتيت بما يكفي ليفقد كلّ أرغونه.

\textbf{20.} ميل الارتفاع بدلالة العمر وفق المربعات الصغرى هو
\qty{20.1}{m.Myr^{-1}}.

\textbf{21.} $(65 - 10)/30 = \qty{1.83}{km}$؛ منذ \qty{43.1}{\Ma}، بمتوسط
$1830/43.1 = \qty{42}{m.Myr^{-1}}$.

\textbf{22.} $\rho\varepsilon/\Lambda = P_0/N - \lambda = \num{3.415e-5} -
\num{0.050e-5} = \qty{3.365e-5}{yr^{-1}}$؛ $\varepsilon = \num{3.365e-5}\times
160/2.7 = \qty{1.99e-3}{cm.yr^{-1}} = \qty{19.9}{m.Myr^{-1}}$.

\textbf{23.} $P_0\Lambda/(\rho N) = \qty{20.2}{m.Myr^{-1}}$، \qty{1.5}{\percent}
أكثر: يكاد التفكّك لا يؤثّر عند هذا المعدّل.

\textbf{24.} $\Lambda/(\rho\varepsilon) = 160/(2.7\times\num{1.99e-3}) =
\qty{30}{kyr}$. ارتفع الغرانيت بنحو \qty{20}{m.Myr^{-1}} بين 52
و\qty{43}{\Ma}، وبمتوسط \qty{42}{m.Myr^{-1}} منذ ذلك، وتتعرّض للتعرية بمعدّل
\qty{20}{m.Myr^{-1}} الآن: كان الانكشاف أسرع في وقت ما
خلال السينوزوي.

\textbf{25.} تبلّر غرانيت فاردن منذ
\textbf{\qty[text-series-to-math=true, separate-uncertainty=true]{311 \pm 10}{\Ma}} (\omterm{def:b3:advanced-geochronology:concordia}{التقاطع العلوي}،
عند انحرافين معياريين).
\end{solution}
